\subsection{Periodic expectations and fixed-support limits}

Write $\Phi_L(A)=V^{(L)}(A)$ for the raw periodic vector and
$Z_L=\langle\Phi_L(A),\Phi_L(A)\rangle$. For a physical matrix $O$ on the
ring, put $\langle O\rangle_L=\langle\Phi_L(A),O\Phi_L(A)\rangle/Z_L$,
with value zero when $Z_L=0$. The trace $\operatorname{Tr}$ below is taken
on the virtual matrix space, whereas $\tr$ is the trace of a virtual matrix.

\begin{theorem}[Finite periodic trace ratio]
    \label{thm:periodic_string_trace_ratio}
    \lean{MPSTensor.mpvExpectation_finKronecker_const_eq_trace_div,
        MPSTensor.physicalObservableTransfer_singleton,
        MPSTensor.physicalObservableTransfer_finKronecker_const}
    \notready
    \uses{def:source_periodic_string_overlap, def:twisted_transfer_map}
    For every tensor $A$, physical matrix $u$, and length $L$, the normalized
    full-ring expectation is
    \begin{align}
        \langle u^{\otimes L}\rangle_L
         & =\frac{\operatorname{Tr}(\E_u^L)}{\operatorname{Tr}(\E_A^L)}.
        \label{eq:periodic_string_trace_ratio}
    \end{align}
    Both sides use the zero convention at a vanishing periodic vector.
\end{theorem}
\begin{proof}
    \uses{lem:ldp_observable_transfer_products}
    The one-site insertion is $\E_u$. Splitting the last physical factor
    and inducting on $L$ gives $\E_{u^{\otimes L}}=\E_u^L$, including $L=0$.
    Expanding the physical numerator and squared norm as endomorphism traces
    gives (\ref{eq:periodic_string_trace_ratio}).
\end{proof}

For the following thermodynamic statements, let $D>0$ and assume
$\E_A(\Id)=\Id$, $\E_A^*(\Lambda)=\Lambda$, $\Lambda>0$, and
$\tr\Lambda=1$. The simple peripheral-space condition means explicitly that
$\E_A(X)=zX$, $X\ne0$, and $|z|=1$ imply $z=1$ and $X\in\C\Id$.
It is the pure canonical regime of
\cite[lines~147--160]{PerezGarcia2008StringOrder}, with the eigenspace
condition retained rather than only uniqueness of the eigenvalue as a set.

\begin{theorem}[Eventual periodic normalization]
    \label{thm:periodic_string_normalization}
    \lean{MPSTensor.pureCanonical_mpvState_inner_self_tendsto_one,
        MPSTensor.pureCanonical_eventually_mpvState_ne_zero}
    \notready
    \uses{def:mpv_inner}
    Under these canonical purity assumptions, $Z_L\to1$ as $L\to\infty$.
    In particular, $\Phi_L(A)\ne0$ for all sufficiently large $L$.
    No nonvanishing assertion at every short length is required.
\end{theorem}
\begin{proof}
    \uses{lem:canonical_transfer_stationary_trace_limit,
        lem:ldp_observable_transfer_products}
    The transfer powers converge pointwise to $P(X)=\tr(\Lambda X)\Id$.
    Finite dimensionality gives convergence of their traces, and the trace
    of the rank-one map $P$ is $\tr\Lambda=1$. Thus
    $Z_L=\operatorname{Tr}(\E_A^L)\to1$.
\end{proof}

\begin{theorem}[Thermodynamic limit with fixed physical support]
    \label{thm:periodic_string_fixed_support}
    \lean{MPSTensor.pureCanonical_mpvExpectation_appendObservable_tendsto,
        MPSTensor.pureCanonical_mpvExpectation_block_string_tendsto,
        MPSTensor.pureCanonical_mpvExpectation_string_tendsto}
    \notready
    \uses{def:physical_string_order_param, def:physical_observable_transfer}
    Under the same canonical purity assumptions, fix a support length $k$
    and a physical operator $O$ on those $k$ sites. Then
    \begin{align}
        \lim_{n\to\infty}\langle O\otimes\Id^{\otimes n}\rangle_{k+n}
         & =\tr\!\left[\Lambda\E_O(\Id)\right].
        \label{eq:periodic_string_fixed_support}
    \end{align}
    In particular, for fixed endpoint blocks $x,y$ of lengths $p,q$ and
    fixed middle length $N$, the limit for
    $x\otimes u^{\otimes N}\otimes y$ is
    $\tr[\Lambda\E_x\E_u^N\E_y(\Id)]$.
    The case $p=q=1$ is display SOPMP of
    \cite[lines~176--181]{PerezGarcia2008StringOrder}.
    The middle length is fixed before the exterior $n$ tends to infinity;
    no exchange of limits or rate uniform in $N$ is asserted.
\end{theorem}
\begin{proof}
    \uses{thm:periodic_string_normalization,
        lem:canonical_transfer_stationary_trace_limit,
        lem:ldp_observable_transfer_products, thm:periodic_string_trace_ratio}
    Put $F=\E_O$ and $P(X)=\tr(\Lambda X)\Id$. The exact numerator is
    $\operatorname{Tr}(F\E_A^n)$. Pointwise convergence and finite
    dimensionality give
    $\operatorname{Tr}(F\E_A^n)\to\operatorname{Tr}(FP)
        =\tr[\Lambda F(\Id)]$.
    The denominator tends to one by
    Theorem~\ref{thm:periodic_string_normalization}. Tensor-product insertions
    give the stated physical string specializations.
\end{proof}

\begin{theorem}[Exact periodic symmetry phase]
    \label{thm:periodic_string_exact_phase}
    \lean{MPSTensor.trace_twistedTransferIter_eq_phase_mul_trace_transfer_pow,
        MPSTensor.mpvExpectation_finKronecker_const_eq_phase_pow}
    \notready
    \uses{def:twisted_transfer_map, def:source_periodic_string_overlap}
    Let $A,u$ be arbitrary and suppose a unitary $V$ and $\mu\in\C$ satisfy
    $\sum_j u_{ij}A_j=\mu V A_iV^\dagger$ for every physical index $i$.
    At every length with $\Phi_L(A)\ne0$,
    \begin{align}
        \langle u^{\otimes L}\rangle_L & =\mu^L.
        \label{eq:periodic_string_exact_phase}
    \end{align}
    The numerator identity
    $\operatorname{Tr}(\E_u^L)=\mu^L\operatorname{Tr}(\E_A^L)$
    holds also at zero periodic vectors, without purity or normalization.
\end{theorem}
\begin{proof}
    \uses{thm:periodic_string_trace_ratio,
        thm:physical_string_phase_retaining_asymptotics}
    The exact iterate identity is
    $\E_u^L=\mu^L L_V\E_A^L L_{V^\dagger}$, where $L_V(X)=VX$.
    Since $L_{V^\dagger}L_V=\Id$, cyclicity of the endomorphism trace
    gives the numerator identity. A nonzero periodic vector has $Z_L\ne0$,
    so (\ref{eq:periodic_string_trace_ratio}) gives
    (\ref{eq:periodic_string_exact_phase}).
\end{proof}

\begin{theorem}[Full-ring spectral alternatives]
    \label{thm:periodic_string_spectral_alternatives}
    \lean{MPSTensor.pureCanonical_mpvExpectation_fullRing_dichotomy,
        MPSTensor.pureCanonical_mpvExpectation_fullRing_tendsto_zero,
        MPSTensor.pureCanonical_mpvExpectation_fullRing_eventually_eq_phase_pow,
        MPSTensor.pureCanonical_mpvExpectation_fullRing_norm_tendsto_one,
        MPSTensor.pureCanonical_mpvExpectation_fullRing_phase_adjusted_tendsto_one,
        MPSTensor.pureCanonical_mpvExpectation_fullRing_not_tendsto}
    \notready
    \uses{def:source_periodic_string_overlap, def:twisted_transfer_map}
    Assume the canonical purity conditions above and let $u$ be unitary.
    Either $\rho(\E_u)<1$ and the complex full-ring expectation tends to
    zero, or $\rho(\E_u)=1$ and its magnitude tends to one.
    In the latter case, choose the unitary intertwiner and phase $\mu$
    supplied by the twisted spectral criterion. The expectation is
    eventually exactly $\mu^L$, so its phase-adjusted value tends to one;
    if $\mu\ne1$, the unadjusted complex sequence has no limit.
    This is the phase qualification of display RL, lines 381--388,
    recorded in~\cite{gap:pgwsvc08_string_order_virtual_boundary}.
    No implication concerning a parameterized Hamiltonian gap is asserted.
\end{theorem}
\begin{proof}
    \uses{lem:pgwsvc08_twisted_spectral_radius,
        thm:periodic_string_trace_ratio, thm:periodic_string_normalization,
        thm:periodic_string_exact_phase}
    Below unit spectral radius, the twisted transfer-power traces tend to
    zero, while $Z_L\to1$. At unit radius the intertwiner gives
    (\ref{eq:periodic_string_exact_phase}) on the eventual nonzero tail.
    If $\mu^L\to z$, then $|z|=1$ and comparison with the one-step shift
    gives $z\mu=z$, forcing $\mu=1$.
\end{proof}
